% !TEX root = ../main.tex
% !TEX program = lualatex
\documentclass[../main.tex]{subfiles}
\IfSubfilesClassLoaded{\externaldocument{\subfix{../build/main}}}{}
% =============================================================================
% APPENDIX: CURRENT EVENTS SNAPSHOT
% =============================================================================

\begin{document}
\IfSubfilesClassLoaded{\chapter{EDO Study Guide}}{}

\chapter{Current Events}\label{app:current-events-snapshot}
\section{Current Events Relevant to EDOs and Acquisition}

\subsection{Project Overmatch and Digital Modernization}
\begin{description}[style=nextline,labelsep=0.5em,font=\bfseries]
  \item[\textbf{Five Eyes Project Arrangement.}] In Oct~2024 Project Overmatch signed a Project Arrangement with Australia, Canada, New Zealand, and the United Kingdom, embedding Cooperative Project Personnel in San Diego to co-develop resilient coalition kill-chain networks ahead of \ac{rimpac}~26; Capt.\ Remil Capili (\ac{pm}), \ac{vadm} Seiko Okano (steering chair/\ac{pmildep}), Shayna Bond (international deputy), and David Flowers (international lead) direct the effort while a coalition lab stands up to support continuous testing~\autocite{DVIDS-Overmatch-FVEY-2025-Nov}.
  \item[\textbf{Open \ac{dagir} pilot.}] By 18~Jun~2025 Project Overmatch, \ac{navwar}, the Chief Digital and Artificial Intelligence Office, and \ac{diu} were onboarding commercial \ac{aiml} applications through the Open \ac{dagir} pipeline so Fleet commanders could visualize readiness data (``Maven Smart Systems'') across the hybrid fleet in weeks instead of years~\autocite{DVIDS-Overmatch-OpenDAGIR-2025}.
  \item[\textbf{Leadership signal.}] \ac{vadm} Seiko Okano continues to serve as the direct reporting program manager and steering committee chair while serving as \ac{pmildep} to \ac{asnrdanda}, keeping Overmatch digital architecture tied to Program Decision Meetings and coalition governance~\autocite{DVIDS-Overmatch-FVEY-2025-Nov,SECNAV5000-2G}.
\end{description}

\subsection{EDO Community Updates}
\begin{description}[style=nextline,labelsep=0.5em,font=\bfseries]
  \item[\textbf{\ac{navwar} change of command.}] Rear Adm.\ (now \ac{vadm}) Seiko Okano relieved Rear Adm.\ Doug Small on 9~Aug~2024, assuming command of \ac{navwar} and direct oversight of Project Overmatch; she emphasized information warfare dominance and continued alliance engagement with industry and academia~\autocite{DVIDS-Okano-2024}.
  \item[\textbf{Role of the \ac{pmildep}.}] As \ac{pmildep} to \ac{asnrdanda}, Okano now chairs Program Decision Meetings when delegated and synchronizes \acp{pae} and \ac{syscom} technical authority, a critical linkage as Capability Trade Councils and portfolio scorecards begin replacing \ac{jcids} artifacts~\autocite{SECNAV5000-2G,secwar-acq-transformation-2025}.
  \item[\textbf{Tycom maintenance focus.}] \ac{usff} and \ac{pacflt} maintenance flag officers continue to drive execution of regional maintenance center reforms and the Naval Sustainment System--Shipyards initiatives highlighted in the 2025 coursebook; these billets remain key board discussion topics because they link \ac{edo} leadership to depot performance~\autocite{edo-3-3-1-acq-policy-players-2025}.
\end{description}

\subsection{Acquisition Policy and Industrial Base Highlights}
\begin{description}[style=nextline,labelsep=0.5em,font=\bfseries]
  \item[\textbf{Authority stack for the \ac{rfo}.}] FAR-overhaul policy direction begins with the presidential action, ``Restoring Common Sense to Federal Procurement''~\autocite{WH-EO-Restoring-Common-Sense-Procurement-2025}. OMB Memo M-25-26 then provides implementation direction to agencies and the \ac{far} Council~\autocite{OMB-M-25-26-Overhauling-FAR-2025}. FAR Council deviation guidance supports interim execution while final rulemaking progresses~\autocite{FAR-Council-Deviation-Guidance-RFO-2025}.
  \item[\textbf{Live implementation tracker.}] Acquisition.gov maintains the live part-by-part deviation tracker for the \ac{rfo}; use it as the operational status source for released text and agency adoption~\autocite{FAR-Overhaul-Deviation-Guide-2026,FAR-Overhaul-RFO-2026}.
  \item[\textbf{Context vs.\ controlling references.}] The White House fact sheet and DAU RFO announcements provide useful explanatory and timeline context, but controlling implementation references remain the presidential action, OMB memo, and FAR Council guidance~\autocite{WH-FactSheet-Restoring-Common-Sense-Procurement-2025,DAU-FAR-Overhaul-Announcement-2025}. Where needed, the Federal Register publication can be used as a metadata anchor for the presidential action record~\autocite{FR-EO-Restoring-Common-Sense-Procurement-2025}.
  \item[\textbf{Current DoD execution.}] DoD has published \ac{rfo} class deviations for \ac{far} Parts~6 (competition), 10 (market research), 12 (commercial products and services), and 18 (emergency acquisitions), all effective Feb~1,~2026; these deviations update \ac{dfars} crosswalks and solicitation prescriptions and should be reflected in acquisition strategies and templates now~\autocite{DoD-RFO-Deviation-Part6-2026,DoD-RFO-Deviation-Part10-2026,DoD-RFO-Deviation-Part12-2026,DoD-RFO-Deviation-Part18-2026}.
  \item[\textbf{\ac{rfo} practical effects.}] The \ac{rfo} is shifting working practice before final rulemaking: rewritten parts arrive first as model deviation text, agencies adopt them through class deviations, and useful but non-mandatory material moves into Strategic Acquisition Guidance, buying guides, and practitioner resources~\autocite{FAR-Overhaul-FAQ-2026,FAR-Overhaul-Deviation-Guide-2026}. Board answer: acquisition teams must check both the codified \ac{far}/\ac{dfars} and the active deviation/adoption status before reusing old templates.
  \item[\textbf{Threshold changes.}] Effective Oct~1,~2025, Acquisition.gov lists inflation adjustments that raise the major-system threshold from \$2.5M to \$3.0M, the micro-purchase threshold from \$10K to \$15K, and the \ac{sat} from \$250K to \$350K~\autocite{AcquisitionGov-Threshold-Changes-2025}. The \ac{fy}~2026 \ac{ndaa} separately raises the certified cost-or-pricing-data threshold for defense contracts entered into after Jun~30,~2026, from \$2M to \$10M and directs \ac{cas} threshold increases, including full \ac{cas} coverage from \$50M to \$100M~\autocite{FY2026-NDAA}.
  \item[\textbf{\ac{acat}/\ac{mdap} limits.}] Section~1804 of the \ac{fy}~2026 \ac{ndaa} resets Title~10 major-system and \ac{mdap} dollar thresholds into \ac{fy}~2024 constant dollars: major-system RDT\&E/procurement thresholds move to \$275M/\$1.3B, and \ac{mdap} RDT\&E/procurement thresholds move to \$1.0B/\$4.5B~\autocite{FY2026-NDAA}. Do not confuse these statutory \ac{acat}-screening thresholds with the \ac{far} major-system threshold used for contracting clauses and procedures.
  \item[\textbf{\ac{was}.}] The Acquisition Transformation Strategy replaces \ac{jcids} with Capability Portfolio Management and Capability Trade Councils, empowering \acp{pae} to adjudicate cost/schedule/performance trades in real time and document decisions in portfolio scorecards~\autocite{secwar-acq-transformation-2025}.
  \item[\textbf{New Navy \acp{pae}.}] On 16~Mar~2026 the \ac{don} announced five additional \ac{pae} organizations---Industrial Operations, Marine Corps, Maritime, Strategic Systems Programs, and Undersea---after the Dec~2025 stand-up of \ac{pae} Robotics and Autonomous Systems~\autocite{DON-PAE-Establishment-2026}. The board-level shift is single-accountable portfolio ownership with authority over program offices plus associated technical, contracting, sustainment, and industrial-base functions.
  \item[\textbf{Industrial base actions.}] \ac{secdef} directed the Department to expand multi-year procurements, establish an Industrial Base Consortium, and reform bid protest incentives so the Munitions Acceleration Council and similar bodies can stabilize demand and attract private capital~\autocite{secwar-acq-transformation-2025}.
  \item[\textbf{\ac{ppbe} integration.}] \ac{usdc} and \ac{usdas} are piloting portfolio budgets and flexible execution authorities so \acp{pae} can move funds inside approved controls while maintaining transparency with \ac{omb} and Congress---a major shift boards will expect candidates to explain~\autocite{secwar-acq-transformation-2025}.
  \item[\textbf{Security cooperation realignment.}] \ac{dsca} and \ac{dtsa} now fall under \ac{usdas}, aligning exportability, Foreign Military Sales, and industrial-base planning so allied orders and U.S. production trades can be brokered inside Capability Trade Councils~\autocite{secwar-arms-transfer-2025}.
  \item[\textbf{National Security Capital Forum.}] Section~1092 of the \ac{fy}~2025 \ac{ndaa} created the \ac{nscf}, chaired by the Office of Strategic Capital, to coordinate public/private financing decisions that affect defense materiel transactions---\acp{edo} should highlight how their programs leverage or are scrutinized by the forum~\autocite{secwar-nscf-2025}.
\end{description}

\subsection{Current Navy Acquisition Issues}
\begin{description}[style=nextline,labelsep=0.5em,font=\bfseries]
  \item[\textbf{2026 \ac{nds} release.}] The latest \ac{nds}, released 23~Jan~2026, replaces the 2022 strategy as the current defense-strategy baseline and aligns to the Nov~2025 \ac{nss}~\autocite{NDS-2026,NSS-2025}. Its four headline lines of effort are homeland defense, deterring China in the Indo-Pacific, increasing allied burden-sharing, and supercharging the U.S. defense industrial base; for \acp{edo}, the acquisition translation is faster fielding, stronger organic and private industrial capacity, and explicit Indo-Pacific sustainment risk management~\autocite{NDS-2026}.
  \item[\textbf{Golden Fleet and \ac{fy}27 budget signal.}] The \ac{don} \ac{fy}~2027 budget request frames the Golden Fleet initiative as a strategy-driven budget, emphasizing shipbuilding, industrial-base revitalization, \acp{pae}, the Rapid Capabilities Office, and disciplined requirements~\autocite{DON-FY27-Budget-Request-2026}. The public brief lists \$377.5B for \ac{don}, \$150B for procurement, 34 total ships, 123 aircraft, \$17B for ship maintenance, and \$1.8B for public shipyards~\autocite{DON-FY27-Budget-Request-2026}.
  \item[\textbf{Make cybersecurity part of our DNA.}] The Department of the Navy Cyber Strategy frames cybersecurity as a warfighting requirement and calls for a cyber-ready culture across the enterprise, reinforcing that programs must treat cyber risk as mission risk~\autocite{DON-CyberStrategy-2023}. DoDI 5000.90 directs program managers to include and decompose cybersecurity requirements from early system development through the full acquisition and sustainment life cycle, including supply chain risk considerations~\autocite{DoDI5000-90}.
  \item[\textbf{Legacy NAVPLAN cues.}] NAVPLAN 2022 and Force Design 2045 remain useful for board context on the hybrid-fleet ambition and Get Real, Get Better culture, but use the 2026 \ac{nds}, \ac{don} \ac{fy}27 budget request, \ac{was}, and \ac{pae} stand-up as the current-event baseline~\autocite{NAVPLAN2022,NDS-2026,DON-FY27-Budget-Request-2026,DON-PAE-Establishment-2026}.
\end{description}
\noindent\textit{Current as of 2026-04-27.}

\section{Shipbuilding and Sustainment Updates}

\subsection{Columbia-Class (SSBN) Delays}
The Columbia-class ballistic missile submarine program, the Navy's highest acquisition priority, experienced significant schedule pressure in 2025. Navy officials reported an estimated 12- to 17-month delay in the delivery of the lead boat, \ac{uss} \textit{District of Columbia} (\ac{ssbn}-826), pushing its first planned patrol into 2030~\autocite{crs-columbia-2025}. The delays were attributed to a combination of factors, including shipyard workforce shortages, supply chain disruptions for critical components like the bow dome and main turbine generators, and the inherent complexities of a first-of-class build~\autocite{crs-columbia-2025,defensenews-columbia-2025}. By late 2025, General Dynamics Electric Boat reported the lead vessel was approximately 60\% complete, but the tight schedule margins remain a significant concern for maintaining the nation's continuous strategic deterrent capability~\autocite{usni-columbia-2025}.

\subsection{Amphibious Warship Block Buy (LPD-17)}
In a significant move to stabilize the amphibious ship industrial base, the Navy executed a multi-ship ``block buy'' procurement for four amphibious warships from \ac{hii} Ingalls in late 2024 and early 2025~\autocite{breakingdefense-amphib-2024}. This procurement, valued at nearly \$1 billion in savings, includes three San Antonio-class amphibious transport docks (\ac{lpd}-33, 34, and 35) and one America-class amphibious assault ship (\ac{lha}-10) to be built between \ac{fy}25 and \ac{fy}29~\autocite{breakingdefense-amphib-2024}. The block buy provides workload stability for the shipyard and its suppliers, enabling more efficient production and ensuring the Navy meets the legally mandated 31-ship amphibious force structure~\autocite{gao-amphib-2024}. This decision followed congressional pressure and signals a commitment to the amphibious fleet after a period of debate over future requirements.

\subsection{Shipyard Infrastructure Optimization Program (SIOP) Status}
The Shipyard Infrastructure Optimization Program (\ac{siop}), a 20-year, \$21 billion effort to modernize the Navy's four public shipyards, reached a critical phase in 2025. A key milestone was the award of a \$377.7 million contract for seismic upgrades to Dry Dock \#4 at Puget Sound Naval Shipyard, which is essential for servicing nuclear submarines, including the future Columbia-class \acp{ssbn}~\autocite{kitsapsun-siop-2025}. However, the program faces significant cost growth due to inflation and unforeseen construction complexities, with annual construction cost escalation exceeding 6\%~\autocite{breakingdefense-siop-2025}. In July 2025, program officials announced a comprehensive review to revise cost and schedule estimates for the entire \ac{siop} portfolio~\autocite{breakingdefense-siop-2025}. Despite these challenges, the Navy continues to invest approximately \$2.7 billion annually, with the goal of reducing submarine maintenance periods by up to 15\% and returning ships to the fleet faster~\autocite{breakingdefense-siop-2025}.

\subsection{West Coast and Pacific Shipyard Development}
West Coast and Pacific infrastructure are now explicit current-event examples of the Foundry/industrial-base story. At Puget Sound Naval Shipyard \& Intermediate Maintenance Facility, a new 25-ton portal crane arrived in July~2025 as part of a \$67M contract for four cranes supporting maintenance, repair, modernization, and retirement of U.S. naval vessels~\autocite{NAVSEA-PSNS-Cranes-2025}. In Hawaii, the future Pearl Harbor Dry Dock~5 remains a major \ac{siop} effort to modernize nuclear-powered submarine maintenance capacity for the Pacific Fleet~\autocite{NAVFAC-PHNSY-DryDock5-CNO-2025,NAVFAC-PHNSY-DryDock5-2023}. In Guam, NAVFAC Pacific awarded a combined \$15B design-build multiple-award construction contract in Sept~2025 for Guam and other Pacific projects; the official release ties the award to resilient, modern, strategically important facilities supporting operational readiness and regional stability~\autocite{NAVFAC-Pacific-Guam-DBMACC-2025}. Board answer: shipyard development is no longer only a facilities topic; it is a naval-operations enabler tied to Indo-Pacific posture, \ac{pae} Industrial Operations, and \ac{nds} industrial-base priorities~\autocite{NDS-2026,DON-PAE-Establishment-2026}.

\IfSubfilesClassLoaded{}{}
\end{document}
